\documentclass[11pt]{article}

\usepackage[T1]{fontenc}
\usepackage{lmodern}
\usepackage{amsmath,amssymb,amsthm,mathtools}
\usepackage[hidelinks]{hyperref}

\newtheoremstyle{harduxdefinition}
{8pt plus 2pt minus 1pt}{8pt plus 2pt minus 1pt}
{\normalfont}{}{\bfseries}{.}{0.5em}{}
\newtheoremstyle{harduxplain}
{8pt plus 2pt minus 1pt}{8pt plus 2pt minus 1pt}
{\itshape}{}{\bfseries}{.}{0.5em}{}
\newtheoremstyle{harduxremark}
{8pt plus 2pt minus 1pt}{8pt plus 2pt minus 1pt}
{\normalfont}{}{\itshape}{.}{0.5em}{}
% \theoremstyle{harduxdefinition}
\newtheorem{definition}{Definition}[section]
\usepackage{amsthm}
% \usepackage{lineno}
% \newtheoremstyle{mydefinition}
%   {22pt}   % 上方空间
%   {22pt}   % 下方空间
%   {\normalfont} % 正文字体
%   {}       % 缩进
%   {\bfseries} % 标题字体
%   {.}      % 标题后标点
%   {0.5em}  % 标题和正文距离
%   {}
% \theoremstyle{mydefinition}
\newtheorem{lemma}{Lemma}[section]
\newtheorem{proposition}{Proposition}[section]
\newtheorem{theorem}{Theorem}[section]
\newtheorem{corollary}{Corollary}[section]
\theoremstyle{harduxremark}
\newtheorem{remark}{Remark}[section]
\newtheorem{example}{Example}[section]

\makeatletter
\def\bstctlcite{\@ifnextchar[{\@bstctlcite}{\@bstctlcite[@auxout]}}
\def\@bstctlcite[#1]#2{\@bsphack
  \@for\@citeb:=#2\do{%
    \edef\@citeb{\expandafter\@firstofone\@citeb}%
    \if@filesw
      \immediate\write\csname #1\endcsname{\string\citation{\@citeb}}%
    \fi}%
  \@esphack}
\makeatother


\DeclareMathOperator{\gra}{gra}

\title{}
\author{Weifeng Yang\\
% Yunnan Earthquake Agency, Kunming, China\\
\texttt{ywf841673182@gmail.com}
}
\date{}

\begin{document}
\bstctlcite{IEEEexample:BSTcontrol}
\linenumbers
\maketitle





\begin{abstract}





\end{abstract}

\noindent\textbf{Keywords:} Haraux function; maximally monotone operator; type~(NI); graph distance; duality mapping.

\medskip
\noindent\textbf{2020 Mathematics Subject Classification:} 47H05; 47N10, 46A20.

\section{Introduction}
\label{sec:introduction}








\subsection{Contributions}

The contributions of this paper are as follows.

(1) 



(2) 



(3) 




The paper is organized as follows. Section~\ref{sec:preliminaries} introduces some preliminary notation and background. Section~\ref{sec:bound} proves the weighted form of quasidensity, establishes the exact decomposition, derives its consequences for the residual infimum, and proves the sharp lower bound. Section~\ref{sec:certificates} investigates exact and approximate residual vanishing, distinguishes the vanishing of the residual infimum from its attainment, and derives a graph distance lower bound for the Fenchel-Young gap.  Conclusions are presented in Section \ref{sec:conclusion}. 






\section{Preliminaries}
\label{sec:preliminaries}









\section{Conclusion}
\label{sec:conclusion}





% \section*{Data Availability Statement}

% No datasets were generated or analyzed during the current study.


% \section*{Funding}

% This work was supported by the CEA Youth Key Project on Earthquake Information (No.~CEAITNS202607), the Spark Program of Earthquake Sciences of China Earthquake Administration (XH25033YB), and the Earthquake Science Technology Innovation Team Project of Yunnan Province (CXTD202507). 

% \section*{Ethics Approval}

% Not applicable.

% \section*{Clinical Trial}

% Not applicable.

% \section*{Competing Interests}

% The author declares that there are no competing interests.

\bibliographystyle{IEEEtran}
\bibliography{refer}

\end{document}
